\documentclass[../main.tex]{subfiles}

\begin{document}
\section{A short proof that $\Q\left[ x \right] / \left( f_n \right)$ has no totally real subfields}
\begin{lemma}
  If $f \in \Q\left[ x \right]$ has a unique real root that then $\Q\left[ x \right] / \left( f \right)$ has no nontrivial totally real subfield.
\end{lemma}
\begin{proof}
  Consider $\overline{\Q}$ an algebraic closure of $\Q$ as a subfield of $\C$, let $G_\Q$ be the total Galois group $\Gal\left( \overline{\Q} / \Q \right)$ and let $\sigma \in G_\Q$ be the element induced by the complex conjugation on $\C$.
  Assume towards contradiction that $\iota: K \hookrightarrow \Q\left[ x \right] / \left( f \right)$ is an embedding of a totally real field $K \neq \Q$.
  Then the restriction induced restriction map 
  \[
    \iota^*: \Hom_\Q\left( \Q\left[ x \right] / \left( f \right), \overline{\Q} \right) \to \Hom_\Q\left( K, \overline{\Q} \right)
  \]
  is a surjective $G$-equivariant map.
  The totally realsness of $K$ implies that $\sigma$ acts trivially on $\Hom_\Q\left( K, \overline{\Q} \right)$, and hence, the $G$-equivariance of $\iota^*$ implies that $\sigma$ reduces to an action on each fiber of $\iota^*$.
  However, the fact that $f$ has a unique real root implies that it has an odd degree, and hence, each fiber of $\iota^*$ has an odd cardinality.
  In particular, since $\sigma$ is an involution, it must have at least one fixed point in each fiber, which contradicts the fact that $f$ has a unique real root.  
\end{proof}
\begin{corollary}
  For each $n$, the quotient $\Q\left[ x \right] / \left( x^n + x^{n-1} + \dots + x - 1 \right)$ has no totally real subfields other than $\Q$.
\end{corollary}
\subsection{Generalization:}
I will try to generalize it to the same generalization as in \cite[Appendix A]{Soroker_appendix_2017}
\begin{lemma}
  Let $E \leq, F, K \leq \overline{E}$ be field extensions such that $E \leq F$ is finite.
  Assume that there exists a prime $p$ such that $p \nmid \left[ F : E \right]$ and such that $\Aut\left( \overline{E}, K \right)$ is a pro-$p$ group.
  Then for any totally $K$ sub-extension $E \leq L \leq F$, 
  \[
    \left[ L : E \right]_s \leq \left | \Hom_E\left( F, K^\ins \right) \right |
  \]
\end{lemma}
\begin{proof}
  Let $G$ be the Galois group $\Gal\left( \overline{E}/ E \right)$.
  Let $\iota: L \hookrightarrow F$ be an $E$-embedding of a totally $K$ extension.
  Consider the restriction map
  \[
    \iota^*: \Hom_E\left( F, \overline{E} \right) \to \Hom_E\left( L, \overline{E} \right)
  \]
  Since $L$ is totally $K$, $\Aut\left( \overline{E} / K \right)$ acts trivially on $\Hom_E\left( L, \overline{E} \right)$.
  Hence, the $G$-equivariance of $\iota^*$ implies that $G$ acts (continuously) on each fiber of $\iota^*$, and in particular, so does the pro-$p$ group $\Aut\left( \overline{E} / K \right)$.
  However, the cardinality of each fiber of $\iota^*$ divides $\left[ F : E \right]$ and hence is prime to $p$.
  Therefore, each fiber of $\iota^*$ has a fixed point under the action of $\Aut\left( \overline{E} / K \right)$, i.e. an element of $\Hom_E\left( F, K^\ins \right)$.
  Therefore, 
  \[
    \left[ L : E \right]_s = \left | \Hom_E\left( L, \overline{E} \right) \right | \leq \left | \Hom_E\left( F, K^\ins \right) \right |
  \]
  as required.
\end{proof}
\begin{remark}
  This is in fact a simplification of the main part of the proof of \cite[Theorem A.5]{Soroker_appendix_2017}.
  So I guess if we copy it and paste instead of that part, the entire appendix can be slightly shortened.
\end{remark}
% We can also prove the following statement which is somewhat similar to \cite[Theorem A.5/A.1]{Soroker_appendix_2017}, but I can't decide yet which one is stronger. 
% \begin{thm}
%   Let $E \leq, F, K \leq \overline{E}$ be field extensions such that $E \leq F$ is finite.
%   Assume that there exists a prime $p$ such that $p \nmid \left[ F : E \right]$ and such that $\Aut\left( \overline{E}, K \right)$ is a pro-$p$ group.
%   Then for any totally $K$ sub-extension $E \leq L \leq F$, 
%   \[
%     \left[ L : E \right]_s \leq \left | \Hom_E\left( F, K^\ins \right) \right |
%   \]
% \end{thm}
% \begin{proof}
  
%   Let $G$ be the Galois group $\Gal\left( \overline{E}/ E \right)$.
%   Let $\iota: L \hookrightarrow F$ be an $E$-embedding of a totally $K$ extension.
%   Consider the restriction map
%   \[
%     \iota^*: \Hom_E\left( F, \overline{E} \right) \to \Hom_E\left( L, \overline{E} \right)
%   \]
%   Since $L$ is totally $K$, $\Aut\left( \overline{E} / K \right)$ acts trivially on $\Hom_E\left( L, \overline{E} \right)$.
%   Hence, the $G$-equivariance of $\iota^*$ implies that $G$ acts (continuously) on each fiber of $\iota^*$, and in particular, so does the pro-$p$ group $\Aut\left( \overline{E} / K \right)$.
%   However, the divisor of each fiber of $\iota^*$ divides $\left[ F : E \right]$ and hence is prime to $p$.
%   Therefore, each fiber of $\iota^*$ has a fixed point under the action of $\Aut\left( \overline{E} / K \right)$, i.e. an element of $\Hom_E\left( F, K^\ins \right)$.
%   Therefore, 
%   \[
%     \left[ L : E \right]_s = \left | \Hom_E\left( L, \overline{E} \right) \right | \leq \left | \Hom_E\left( F, K^\ins \right) \right |
%   \]
%   as required.
% \end{proof}



\end{document}